
\section{Introduction}
\label{sec:introduction}

Concurrency is no longer an implementation detail. To make effective
use of multi-core hardware and distributed infrastructure, modern
software systems decompose their work into cooperating activities that
may run in different threads, on different cores, or on different
machines. These activities cannot remain isolated: they must exchange
data, coordinate progress, and agree on the order of their
interactions. Getting these protocols right is difficult. A missed
message, an unexpected choice, or a communication step performed in
the wrong state can turn an otherwise natural decomposition into a
source of subtle failures. A robust language discipline should
therefore specify both the data exchanged by concurrent components
and the order in which those exchanges may occur.

% TODO [Importance: High; Urgency: High]: Replace the current placeholders with
% a complete problem statement: session types ensure protocol fidelity, CFSTs
% make non-tail-recursive protocols expressible, but programming with them still
% forces resource-passing style and polymorphic recursion.
% TODO [Importance: High; Urgency: High]: Define borrowing early and concretely:
% a borrow consumes a prefix of a session and leaves a residual suffix, with the
% typing discipline internalizing the linear update of the channel.
% TODO [Importance: High; Urgency: High]: Explain why borrowing for CFSTs is
% technically different from borrowing for regular session types: sequential
% composition creates non-tail protocol prefixes, so splitting must work modulo
% CFST equality and across choices/recursion.
% TODO [Importance: High; Urgency: High]: Preview the local-vs-remote borrowing
% distinction from the motivation section, including why local borrows need no
% synchronization and remote borrows introduce Ret/Acq handover points.
% TODO [Importance: High; Urgency: Medium]: Turn the comparison with BGV into a
% precise positioning paragraph: regular rather than context-free session
% types, semantics by translation rather than direct operational semantics, and
% uniform remote-style treatment of borrows.
% TODO [Importance: High; Urgency: High]: Align the contribution list with the
% body: direct source semantics, declarative type system with BI-contexts,
% algorithmic typing/constraint solving, target calculus, translation
% simulation or bisimulation, implementation, and mechanization status.
% TODO [Importance: Medium; Urgency: High]: State the current status of the
% metatheory accurately. In particular, distinguish intended type soundness,
% algorithmic soundness/completeness, and translation simulation/bisimulation
% from results that are already proved or mechanized.
% TODO [Importance: Medium; Urgency: Medium]: Introduce the running examples
% used later, especially tree serialization for local borrows and HTML rendering
% for remote borrows, so the motivation section does not have to establish all
% context from scratch.
% TODO [Importance: Medium; Urgency: Medium]: Add a short paper-organization
% paragraph once the final section structure is stable.
% TODO [Importance: Low; Urgency: Medium]: Normalize terminology in the
% introduction: \thecalculus, context-free session types with borrowing,
% resource-passing style, local borrow, remote borrow, Ret, and Acq should be
% introduced before they are used in the contribution bullets.


Session types
\cite{DBLP:conf/concur/Honda93,HondaVK98,Gay_Vasconcelos_2025}
address precisely this problem: they are a type discipline for
structured bidirectional communication
protocols. Their core features are typed, bidirectional message
exchange and typed choices. They ensure protocol fidelity as well as type soundness,
when embedded in concurrent lambda calculus. There is significant
literature on session types, which is well appreciated in some recent 
surveys and books
\cite{DBLP:journals/csur/HuttelLVCCDMPRT16,
  gay17:_behav_types,
  DBLP:journals/ftpl/AnconaBB0CDGGGH16,
  Gay_Vasconcelos_2025}. 
While much recent work concentrates on multi-party session types
\cite{DBLP:journals/jacm/HondaYC16}, this work is solely concerned
with binary session types, which govern two-party protocols.

Session types support infinite and unbounded behaviors through
recursive types almost from the beginning \cite{HondaVK98}. Since
then, most work has concentrated on \emph{regular} recursive session
types, where recursion is restricted to tail positions in the
type. For example, here is the type
\lstinline|mualpha.&{next:?Int.alpha, quit:End}|
of a server that receives a list of integers, where the occurrence of
the recursive \lstinline{alpha} is syntactically restricted to tail
position. Context-free session types
\cite{DBLP:conf/icfp/ThiemannV16,DBLP:journals/iandc/AlmeidaMTV22} (CFST)
have lifted this restriction by adopting a symmetric composition
operator \lstinline{_;_} rather than prefixing
\lstinline{?Int._} as in regular session types. Symmetric composition
enables more expressive communication structures, like trees and
stacks, but it requires a non-trivial notion of type conversion.

Borrowing has received much attention recently due to its adoption in
the Rust language \cite{rust-reference}. It refers to methods that regulate the
ownership transfer of resources between different parts of a
program. Originally conceived for imperative programming,
\citet{DBLP:journals/pacmpl/SaffrichS0V25} present a session type
system with a notion of borrowing.
In this system, a borrow consumes a prefix of a session and leaves a residual suffix, with the
typing discipline internalizing the linear update of the channel.
Their calculus BGV is based on regular
session types \cite{DBLP:conf/esop/LindleyM15} and its semantics is
given by a translation to a calculus from the literature
\cite{DBLP:journals/lmcs/KokkeD23}. The approach of the BGV paper has a number of
drawbacks. As the source calculus has no direct semantics, it is
difficult to perform optimizations at the source level. Moreover, the
translation is uniform in the sense that it does not distinguish
between local borrows that remain in the same process and remote
borrows that are passed to other processes. This uniformity leads to
inefficiency because local borrows can elide the synchronization
machinery needed for remote borrows. 



\subsection*{Contributions}

\begin{itemize}
\item A new design of session types with borrowing, simpler than previous work.
\item Based on context-free session types, rather than regular session types.
\item Splitting for context-free session types.
\end{itemize}

Technical contributions
\begin{itemize}
\item Source language with direct operational semantics, rather than translational. The
  operational semantics serves as an abstract specification of
  borrowing for session types.
\item New intermediate language that makes channel handling explicit
  and efficient.
\item Translation from source language to intermediate language: maps
  specification to implementation.
\item Source language: Type system with effects based on BI-contexts.
\item Metatheory: type soundness wrt.\ operational semantics,
  simulation for the translation.
\item Implementation: type checker leveraging the implementation of
  FreeST \cite{almeida26:_frees_concur_progr_languag_with,DBLP:journals/corr/abs-1904-01284}
\end{itemize}

%%% Local Variables:
%%% mode: latex
%%% TeX-master: "../popl2027.tex"
%%% End:
